%=============================================================================
% Section 6: The Action Interface Definition
%=============================================================================
\section{The Action Interface Definition}

\subsection{Formal Definition}

We define the action interface formally:

\begin{definition}
The Action Interface is the boundary between AI cognition and the capability ecosystem, defined by:
\[
\text{ActionInterface} = (\text{Capabilities}, \text{Routing}, \text{Execution})
\]
Where:
\begin{itemize}[leftmargin=*]
\item \textbf{Capabilities:} The set of available capabilities $C = \{c_1, c_2, \ldots, c_n\}$
\item \textbf{Routing:} A function $\text{Route}(c, \text{context}) \to \text{provider}$
\item \textbf{Execution:} A function $\text{Execute}(c, \text{input}, \text{provider}) \to \text{output}$
\end{itemize}
\end{definition}

\textbf{PolicyDecision Formalization:} The routing decision process can be formalized as:
\[
\text{PolicyDecision}(c, \text{ctx}) \to (\text{provider}, \text{constraints}, \text{metadata})
\]

Where the input includes $c$ (the capability being requested) and $\text{ctx}$ (execution context containing preferences, constraints, and history), and the output includes the selected provider, applied constraints, and decision rationale.

\textbf{Decision Logic:}
\[
\text{PolicyDecision}(c, \text{ctx}) = \arg\min_{p \in \text{Providers}(c)} \text{Score}(p, \text{pref}, \text{cons})
\]

Where:
\[
\text{Score}(p, \text{pref}, \text{cons}) = w_1 \cdot \text{LatencyScore}(p) + w_2 \cdot \text{CostScore}(p) + w_3 \cdot \text{QualityScore}(p) + \text{Penalty}(p, \text{cons})
\]

with $w_1 + w_2 + w_3 = 1.0$ and $\text{Penalty}(p, \text{cons}) = 0$ if $p$ satisfies all constraints, $\infty$ otherwise.

\textbf{Dynamic Weight Override:} The preference weights $(w_1, w_2, w_3)$ support dynamic override by the Cognition Layer at request time. This enables agents to adjust routing behavior based on task urgency or context.

This formalization enables transparency (decisions can be explained), optimization (scoring functions can be tuned), and compliance (constraints are explicitly represented).

\subsection{Precision Statement}

We can state the role of the contact layer precisely:

\begin{quote}
\textit{The Contact Layer is the policy and routing component of the action interface between AI cognition and the capability ecosystem, responsible for capability routing, provider selection, and policy evaluation.}
\end{quote}

This definition captures:
\begin{itemize}[leftmargin=*]
\item \textbf{Position:} Between cognition and execution runtime
\item \textbf{Function:} Routing, policy evaluation, provider selection
\item \textbf{Scope:} Control plane of the Action Layer
\item \textbf{Domain:} Capability mediation (not execution mechanics)
\end{itemize}

\textbf{Note on Formalization:} The PolicyDecision formalization above is conceptual, illustrating the decision logic framework. Detailed algorithm implementations, including constraint satisfaction, weight learning, and multi-objective optimization, are addressed in the ai-lib reference implementations.

\subsection{Architectural Significance}

This definition has architectural implications:
\begin{enumerate}[label=\arabic*.]
\item \textbf{Separation of Concerns:} Cognition and action are architecturally distinct, enabling independent evolution.
\item \textbf{Interface Contracts:} The action layer defines clear contracts for capability invocation, enabling provider substitution.
\item \textbf{Ecosystem Participation:} By defining the action interface, systems can participate in capability ecosystems.
\item \textbf{Decoupling Principle:} Agents need not know provider details---they express capability requirements, and the action layer handles implementation.
\end{enumerate}
